\chapter{DISCUSSION}

\section{From Static Maps to Dynamic Visual Networks}
Historically, the functional organization of the primary visual system has been delineated through the static mapping of receptive fields and anatomical inter-hemispheric tracing \cite{wunderle2013}; \cite{schmidt2023}. While these classical methods successfully outline the spatial constraints and functional columns of the cortex, they inherently capture only a time-averaged "structural backbone." Cognitive processing and the integration of moving stimuli, however, depend on millisecond-scale state transitions within neural ensembles—a temporal dimension of rapid network evolution \cite{kopell2014}. 

Even as our anatomical maps reach unprecedented levels of precision---such as recent efforts updating the classic Binzegger cortical wiring diagram to reveal extensive, previously unaccounted local excitatory synapses terminating in layer 6 \cite{Martin_2024}---structure alone cannot predict traffic. While anatomy dictates the available pathways, only time-resolved effective connectivity frameworks can reveal how this intricate wiring operates on a millisecond scale.

The application of the NECTA framework to the \texttt{c1608a01} dataset provides empirical evidence of this dynamic nature. By utilizing a continuous sliding-window approach, our results suggest that the visual cortex undergoes rapid, stimulus-locked topological reorganizations, reaching transient peaks in \textit{Small-Worldness} precisely during the moving grating stimulus.

This macroscopic network efficiency---an active, responsive property rather than a fixed anatomical trait---is strongly supported by recent cellular and anatomical discoveries. On a microscopic scale, \textit{in vitro} microfluidic studies by \cite{Lassers2026} have shown that theta oscillations (4--10 Hz) in isolated axons occur spontaneously and independently of spiking. These oscillations actively coordinate target neuron bursting based on phase and amplitude through a sparse coding architecture, where only a restricted fraction of axons carries high-amplitude signals. This selective, sparse recruitment of communication channels likely mediates the transient optimization detected by NECTA.

Furthermore, \cite{Martin_2024} identified an extraordinarily extensive horizontal spread within Area 17, with neurons integrating information across distances of up to 11 millimeters (spanning 5--10 degrees of the visual field). This vast lateral arborization provides a plausible anatomical substrate for the network to rapidly bind global spatial information, explaining the acute surges in \textit{Small-Worldness} observed computationally. Their discovery of a novel, direct ``shortcut'' projection from the layer 3/4 border to layer 6---bypassing the canonical layer 5 relay and likely associated with the fast-conducting, motion-sensitive Y-pathway---further underscores the necessity of high-frequency dynamic analyses. The existence of such fast parallel anatomical routes suggests that visual processing is not strictly serial or slow, fully justifying our focus on the gamma band and sliding-window techniques to capture rapid synchronizations that static tools completely miss.

\section{Isolating the Feedforward Flow in the Anesthetized State}
A critical strength of our biological validation lies in the physiological state of the animal model. The recordings were obtained under general anesthesia and complete paralysis. Historically, anesthesia has often been viewed as a limitation for studying information flow. However, within the modern framework of predictive coding, this state acts as an analytically significant "pharmacological lesion" model. The literature indicates that general anesthesia preferentially suppresses the inhibitory top-down feedback pathways, which are typically mediated by slower alpha and beta oscillations \cite{vezoli2021}. 

Crucially, while top-down executive modulations are suppressed, the bottom-up feedforward pathways---mediated by high-frequency gamma oscillations---are left to operate autonomously. Therefore, rather than a mere limitation, the anesthetized state elegantly isolated the intrinsic feedforward flow from Area 17 to Area 18. This pharmacological isolation provides a powerful justification for our findings of strong, stimulus-locked directionality restricted to the low-gamma band. It guarantees that the dynamic topological reorganizations and the causal directionality shifts extracted by NECTA are not artifacts of top-down cognitive states, attention, or saccadic eye movements. Instead, the framework successfully unveils the pure, low-level reflexive plasticity inherent to the local synaptic loops and dense transcallosal pathways interlinking A17, A18, and the TZ.

\section{Mitigating Electrophysiological Confounders}
A significant challenge in high-density LFP connectomics is the pervasive artifact of volume conduction, wherein a single electrical source is instantaneously registered by multiple neighboring electrodes \cite{bastos2015}. If unaccounted for, this phenomenon generates spurious zero-phase lag correlations, heavily biasing the resulting graph topology with false-positive edges. 

NECTA addresses this methodological bottleneck architecturally. By employing the Hilbert Transform to extract analytical signals, the framework calculates phase-based indices such as the Weighted Phase Lag Index (\textit{wPLI}) and Imaginary Coherence (\textit{iCoh}), which are mathematically immune to instantaneous spatial spread. This algorithmic refinement facilitates the extraction of dynamic topologies that reflect true underlying physiological synchronization rather than electrical crosstalk. This distinction is critically reinforced by recent experimental evidence isolating single axons in microfluidic tunnels, which suggested that subthreshold voltage oscillations are genuine, regenerative local phenomena of the axon rather than mere volume-conducted epiphenomena of the extracellular matrix \cite{Lassers2026}. By utilizing \textit{wPLI} and \textit{iCoh}, NECTA's architecture is consistent with this biological reality, mathematically isolating true phase-delayed interactions.

\section{Gamma Synchronization, Predictive Coding, and Callosal Modulation}
Previous work from our group has established the cardinal role of gamma-band synchronization in functional routing and cohesion within the retinogeniculate and cortical pathways of the cat \cite{neuenschwander2023}. NECTA extends this foundation by mapping the causal directionality of these high-frequency rhythms. 

By applying Partial Directed Coherence (PDC) optimized dynamically via the Bayesian Information Criterion (BIC), we eliminated the bias of static model-order selection (which often leads to overfitting when estimating temporal shifts, unlike AIC). The resulting hierarchy in the Low Gamma band (A17 acting as the persistent Driver, A18 as the Receiver) is consistent with current theories of predictive coding and feedforward information routing \cite{vezoli2021}. At the cellular level, this macroscopic directionality is supported by findings that the phase of axonal theta oscillations dictates the firing probability of target neurons, establishing a clear physical causality where subthreshold oscillations `prime' target neurons for spiking \cite{Lassers2026}. 

A fundamental hurdle when applying multivariate autoregressive (MVAR) based metrics like PDC to continuous electrophysiological recordings is the assumption of signal stationarity, which is heavily violated during sensory processing, particularly by transient evoked potentials. To robustly address this, NECTA tackles non-stationarity through a rigorous multi-stage approach. First, adopting a classical defense well-established in the Granger causality literature, the pipeline performs the subtraction of the ensemble average (the mean evoked potential) from each individual trial. By removing this phase-locked evoked response, the framework successfully isolates purely induced cerebral activity, effectively mitigating the primary source of signal non-stationarity within short sliding windows (e.g., 500 ms). Second, it implements a stationarity filter based on stringent variance thresholds. This filter actively identifies and isolates any remaining epochs of extreme amplitude volatility, ensuring that the sliding windows processed by the MVAR primarily contain stable oscillatory activity. Finally, for windows exhibiting mild non-stationarity breaks, NECTA relies on the dynamic optimization of the Bayesian Information Criterion (BIC). As demonstrated by \cite{Faes_2010}, estimating causality in the frequency domain requires precise control over model complexity to handle instantaneous and lagged interactions without overfitting. The stringent penalty parameter intrinsic to BIC actively limits the autoregressive model order during these state shifts, effectively penalizing spurious complexity. Consequently, this dynamic scaling prevents the PDC estimates from modeling transient artifacts or noise as causal edges, ensuring that the detected causal hierarchies reflect genuine physiological routing.

The validity of utilizing PDC for this structural mapping is strongly supported by recent large-scale computational studies. Nunes et al. (2021) demonstrated \textit{in silico} using spiking neuronal networks that generalized PDC provides a highly reliable estimate of true underlying anatomical structural connectivity (yielding a Pearson correlation of $r=0.74$). Crucially, their model intrinsically generated LFP signals with prominent spectral peaks in the gamma band, strongly supporting our empirical focus on low-gamma (30--59 Hz) synchronization as a primary channel for causal routing. Furthermore, a persistent challenge in high-density electrophysiology is the ``partial observation'' problem, where unrecorded brain regions might introduce spurious causalities. \cite{Nunes_2021} indicated that PDC estimates maintain robust correlations with structural connectivity ($r > 0.6$) even when the analysis is restricted to a small spatial cluster. This provides powerful theoretical justification that the causal inferences drawn by NECTA from the restricted A17-TZ-A18 microcircuit remain statistically valid despite the absence of whole-brain recordings. At the cellular level, this macroscopic directionality is supported by findings that the phase of axonal theta oscillations dictates the firing probability of target neurons, establishing a clear physical causality where subthreshold oscillations `prime' target neurons for spiking \cite{Lassers2026}. 

Furthermore, identifying the Transition Zone (TZ) as a "Dynamic Modulator" that reverses its causal flow upon stimulus onset highlights the necessity of time-resolved directional tools. Anatomically, the TZ boundary is densely innervated by mixed intra-hemispheric projections and transcallosal fibers that align spatial orientation maps across the vertical meridian \cite{schmidt2023}. By switching from a passive sink to an active causal source during the moving grating stimulation, the TZ likely orchestrates the interhemispheric temporal unifications required to track sensory borders crossing visual fields. Static connectivity would have merely averaged this reversal into an indiscernible null-flow, masking a critical operational mechanism of visual processing. This substantial functional volatility is theoretically supported by the concept of ``activity flow'' across neural networks, where the variability of a node's directed functional connectivity is heavily influenced by its structural centrality \cite{Nunes_2021}. Crucially, this dynamic hub-like behavior is grounded in recent retrograde tracing data, which revealed that the border region representing the vertical meridian (the TZ) acts as an epicenter of convergence for both dense long-range ipsilateral projections and substantial contralateral inputs \cite{Martin_2024}. This anatomical reality explains the results observed in our Node Strength analysis, where the spatially restricted TZ statistically surpassed the substantial extrastriate Area 18 in topological centrality. Despite Area 18 being a large region responsible for processing a substantial volume of complex visual features, the TZ's structural wiring as a dense convergence zone causes it to act as an interhemispheric bottleneck under high connectivity load. Thus, the TZ is anatomically predetermined to act as a highly labile dynamic modulator within the A17 $\rightarrow$ TZ $\rightarrow$ A18 causal network, gating information flow based on the immediate demands of the moving stimulus.


\section{Overcoming the Computational Scale and Open Science}
Translating the theoretical concepts of dynamic connectomics into practical analytical outputs requires overcoming severe computational barriers. Generating stochastic multivariate null models (e.g., surrogate data phase-randomization and configuration models) across continuous sliding windows triggers a combinatorial explosion that typically crashes single-threaded analytical scripts. 
By orchestrating these routines through a Dask-distributed High-Performance Computing (HPC) backend and utilizing `xarray` and `Zarr` for lazy-loading operations, NECTA effectively eliminates the "Big Data" electrophysiology bottleneck. 

Equally vital to this computational leap is the aspiration toward \textit{Open Science} and reproducibility. While systemic barriers remain in widespread adoption, the rigorous parameterization, GitHub/CodeMeta logging, and containerized deployment of NECTA’s codebase represent a significant effort to ensure that our analytical workflows are transparent and globally verifiable. By doing so, it provides the ideal computational environment to validate large-scale biological theories, such as integrating subthreshold oscillatory dynamics into next-generation Spiking Neural Networks (SNNs) for more energy-efficient temporal processing \cite{Lassers2026}. Equally vital to this computational leap is the adherence to \textit{Open Science} and reproducibility. The rigorous parameterization, GitHub/CodeMeta logging, and containerized deployment of NECTA’s codebase ensure that our analytical workflows are transparent and globally verifiable. This pipeline aims to lower the barriers to complex graph-theoretical analysis, striving to allow the scientific community to interactively slice and visualize multidimensional connectomes.

